% ---------- 页面与基础 ----------
\usepackage{geometry} % 页面尺寸与页边距设置
\geometry{margin=2.5cm}
\usepackage{indentfirst} % 使章节后的首段缩进
\setlength{\parindent}{2em}
\usepackage{placeins} % 提供 \FloatBarrier 控制浮动体位置

% ---------- 数学与公式 ----------
\usepackage{amsmath,amssymb} % 高级数学公式环境与数学符号

% ---------- 图形与颜色 ----------
\usepackage[table]{xcolor} % 颜色支持（含表格着色）
\usepackage{graphicx} % 插入图片命令 \includegraphics
\usepackage{tikz} % 矢量绘图，配合 pgfplots 使用
\usetikzlibrary{positioning, calc, shapes.geometric, arrows.meta, backgrounds, external, fit} % TikZ 常用库

% ---------- 表格 ----------
\usepackage{booktabs} % 专业表格线命令（\toprule/\midrule/\bottomrule）
\usepackage{tabularx} % 可伸缩列宽的表格环境
\usepackage{array} % 表格列格式增强
\usepackage{multirow} % 合并多行单元格
\usepackage{longtable} % 跨页长表格
\usepackage{colortbl} % 表格单元格背景色支持

% ---------- 列表、代码与文本 ----------
\usepackage{enumitem} % 列表定制（间距、标签等）
\usepackage{soul} % 文本高亮（\hl）和删除线等
\usepackage{listings} % 源代码高亮与代码块样式
%========================================
% Emoji support using twemojis (Vector SVGs)
%========================================
\usepackage{twemojis} % 将 emoji 渲染为矢量图（twemoji SVG）
\usepackage{alltt} % 类似 verbatim 但允许 LaTeX 命令展开（用于示例输出）
\usepackage{alltt}

\usepackage{subfiles} % 支持将章节作为独立子文件编译
\usepackage[numbers,sort&compress]{natbib} % 参考文献样式与引用压缩

% ---------- 图题与浮动 ----------
\usepackage{float} % 强制浮动体位置（H）参数
\usepackage{caption} % 自定义图表题注样式
\usepackage{subcaption} % 子图题注环境

% ---------- 其他常用工具 ----------
\usepackage{changelog} % 变更日志支持（项目文档用途）
\usepackage{tcolorbox} % 彩色盒子用于提示、示例等
\tcbuselibrary{skins,breakable} % tcolorbox 的扩展库
\usepackage{etoolbox} % 编程型宏工具集合
\usepackage{siunitx} % 物理量与单位格式化

% ---------- TODO 注释 ----------
% 默认使用页边 TODO，尽量不影响正文排版。
\usepackage[textsize=scriptsize]{todonotes} % 文档内待办注释（边注）
\setlength{\marginparwidth}{2.1cm}

% ---- TikZ 外部化缓存 (首次编译生成 PDF，后续直接复用，大幅提速) ----
% 编译时必须开启 -shell-escape（VSCode 配置已设置）
% 如需临时关闭缓存（调试用），注释下面两行并取消 \tikzexternaldisable 的注释
% prefix 只能是项目根目录下的相对路径（不含 out/ 前缀），
% 因为 xelatex 的 -output-directory=out 会让 openout 自动在前面加 out/，
% 若 prefix 也写 out/ 则变成 out/out/tikz-cache/ —— 路径叠加导致写入失败。
% .md5 文件写到 out/tikz-cache/（kpathsea 会把 out/ 加入搜索路径，读取正常）；
% 外部化子进程（无 -output-directory）将 PDF 写到项目根 tikz-cache/，\includegraphics 可直接找到。
\tikzexternalize[prefix=tikz-cache/]
%\tikzexternaldisable  % 取消注释可关闭外部化，回退到内联编译（调试用）

% ---- 修复 tcolorbox (skins+breakable) 与 TikZ external 的兼容性冲突 ----
% tcolorbox skins 库内部会创建 tikzpicture 用于绘制边框/背景；
% TikZ external 全局 hook 会拦截这些内部 tikzpicture，在 breakable 跨页时
% 导致 collect@until@end@tikzpicture 状态错乱（Extra \else / extra } 等连锁错误）。
% 解决方案：在每个 tcolorbox / personalnote 环境的开头/结尾禁用/启用外部化。
\AtBeginEnvironment{tcolorbox}{\tikzexternaldisable}
\AtEndEnvironment{tcolorbox}{\tikzexternalenable}
\AtBeginEnvironment{personalnote}{\tikzexternaldisable}
\AtEndEnvironment{personalnote}{\tikzexternalenable}
\AtBeginEnvironment{taskbox}{\tikzexternaldisable}
\AtEndEnvironment{taskbox}{\tikzexternalenable}

\definecolor{ieeeBlue}{RGB}{0, 83, 159}
\definecolor{ieeeRed}{RGB}{198, 38, 46}
\definecolor{ieeeGreen}{RGB}{0, 114, 41}
\definecolor{ieeeOrange}{RGB}{235, 110, 31}
\definecolor{ieeeGray}{RGB}{100, 100, 100}
\definecolor{darkGray}{RGB}{66, 66, 66}
\definecolor{lightGray}{RGB}{189, 189, 189}

% TODO 全局样式与便捷命令
\setuptodonotes{
  color=ieeeOrange!15,
  bordercolor=ieeeOrange,
  linecolor=ieeeOrange
}
% todonotes 内部依赖 TikZ；开启 externalize 时需临时禁用外部化，避免编译失败。
\let\origTodo\todo
\RenewDocumentCommand{\todo}{ O{} m }{%
  \tikzexternaldisable
  \origTodo[#1]{#2}%
  \tikzexternalenable
}
\newcommand{\todoinline}[1]{\todo[inline,color=ieeeOrange!12]{#1}}

\usepackage{pgfplots} % 基于 TikZ 的函数与数据绘图
\pgfplotsset{compat=1.18}


% ctex 设置，确保各级标题缩进符合中文规范
\ctexset{
  paragraph/indent = 2em,
  subparagraph/indent = 2em,
}



% 链接相关宏包建议放在导言区靠后位置
\usepackage{hyperref}
\hypersetup{
  colorlinks=true,
  linkcolor=blue,
  urlcolor=blue,
  pdftitle={知识库文档},
  pdfauthor={葛良晨}
}

\lstdefinelanguage{riscv-asm}{
  morekeywords={add, addi, sub, sll, slli, srl, srli, sra, srai, and, andi, or, ori, xor, xori, slt, slti, sltu, sltiu,
    beq, bne, blt, bge, bltu, bgeu, jal, jalr, lui, auipc, lw, lh, lb, lhu, lbu, sw, sh, sb,
    csrr, csrw, csrrs, csrrc, csrrwi, csrrsi, csrrci, mret, ecall, ebreak, wfi, nop, ret, call, j, li, la, mv,
    .text, .data, .bss, .rodata, .section, .global, .globl, .align, .file, .type, .size, .word, .half, .byte, .asciz},
  morecomment=[l]{\#},
  morecomment=[l]{;},
  morestring=[b]",
  sensitive=true
}

\lstdefinelanguage{SystemVerilog}{
  morekeywords={
    always_comb, always_ff, always_latch, assert, assume, automatic, begin, end, case, endcase,
    class, endclass, clocking, endclocking, constraint, endconstraint, cover, covergroup, endgroup,
    default, disable, do, else, endfunction, endmodule, endpackage, endprogram, endtask, enum,
    export, extends, final, for, foreach, fork, join, join_any, join_none, function, generate,
    endgenerate, genvar, if, import, inside, interface, endinterface, localparam, logic, modport,
    module, new, package, program, protected, rand, randc, randcase, randsequence, ref, return,
    sequence, endsequence, struct, task, typedef, union, unique, unique0, virtual, void, wait, while,
    with
  },
  morecomment=[l]{//},
  morecomment=[s]{/*}{*/},
  morestring=[b]",
  sensitive=true
}

\lstset{
  basicstyle=\ttfamily\small,
  frame=single,
  breaklines=true,
  keywordstyle=\color{blue},
  commentstyle=\color{teal},
  stringstyle=\color{red!70!black}
}

% ---------- 自定义行内代码样式 ----------
\definecolor{codebg}{RGB}{230, 230, 230}
% 定义新命令 \code，避免直接修改 \texttt 导致系统崩溃
\newcommand{\code}[1]{\colorbox{codebg}{\ttfamily #1}}

% 保持 tcolorbox 的兼容性设置
\tcbset{shield externalize}





\title{RISC-V 技术文档 Demo}
\author{葛良晨}
\date{\today}

% 自定义个人总结环境
\newtcolorbox{personalnote}[1][]{
  colback=yellow!10,
  colframe=yellow!70!black,
  fonttitle=\bfseries,
  title=个人总结,
  breakable,
  enhanced,
  before upper={\setlength{\parindent}{2em}},
  #1
}

% 自定义待办事项环境
\newtcolorbox{taskbox}[1][]{
  colback=blue!5,
  colframe=blue!60!cyan!70!black,
  fonttitle=\bfseries,
  title=待办事项,
  breakable,
  enhanced,
  before upper={\setlength{\parindent}{2em}},
  #1
}

% 自定义妙妙想法彩框
\newtcolorbox{ideabox}[1][]{
  colback=purple!8!white,
  colframe=purple!60!black,
  fonttitle=\bfseries,
  title=✨ 妙妙想法,
  breakable,
  enhanced,
  before upper={\setlength{\parindent}{2em}},
  #1
}


